\documentclass[preview]{standalone}
\usepackage{amsmath}
\usepackage{booktabs}
\usepackage{amssymb}
\usepackage{setspace}
\usepackage{bm}
\let\oldcdot\cdot
\usepackage{breqn}
\let\cdot\oldcdot
\begin{document}
Let $\phi(x)$ be the level-set function and
 \begin{align*}\raisebox{\depth}{$\chi$}(x) &= \begin{cases} 1 & \mbox{if}\quad \phi(x) < 0 \\ k_{air}/k_0 & \mbox{if}\quad \phi(x)  \ge 0 \end{cases},\end{align*}
where the weak-material contrast is the physically motivated ratio of the thermal conductivity of air to that of aluminium, $k_{air}/k_0 = 0.026/200 \approx 1.3 \times 10^{-4}$ (numerically $1.3 \times 10^{-4}$, the value taken by the characteristic function).
The design domain is the annular cross-section of a district-heating pipe, $D = \left\{ \mathbf x : r_0 < \left|\mathbf x\right| < R \right\}$ with outer radius $R = 0.5$ m (a pipe of diameter 1 m, scaled to unit size for benchmark comparison) and initial inner radius $r_0 = 0.1$ m.  The OUTER circle $\Gamma_o$ is a fixed boundary: the pipe casing exchanges heat with the environment at $u_{amb} = 298$ K through a Robin condition with effective convective heat transfer coefficient $h = 25$ W/(m$^{2}$ K).  The INNER circle $\Gamma_i$ is a design boundary — the tube wall whose shape is optimised — swept by the district-heating water at temperature $u_H = 353$ K, which injects heat by convection through a Robin condition with heat transfer coefficient $h_H = 1000$ W/(m$^{2}$ K) — a realistic forced-convection coefficient for the water side of a heating network.  The optimisation maximises the heat rejected to the environment through $\Gamma_o$.
The Lagrangian of the convectively-heated pipe problem is
\begin{dmath*}
\mathcal L(\left[ u\right],\left[ \delta u\right]) = J(u) +  \int_{D} k \raisebox{\depth}{$\chi$} \left(  \nabla u  \cdot  \nabla \delta u  \right) \,dx + \int_{\Gamma_i} h_H \left( u - u_{H}\right) \delta u \,ds_x + \int_{\Gamma_o} h \left( u - u_{amb}\right) \delta u \,ds_x ,
\end{dmath*}
with
\begin{align*}
J(u) &=  \int_{D} \ell \raisebox{\depth}{$\chi$} \,dx + \int_{\Gamma_o} q_b \left( u_{amb} - u\right) \,ds_x ,
 \end{align*}
where the outer-surface temperature term is weighted by $q_b = h_H \left( u_H - u_{amb}\right) = 1000 \times 55 = 5.5 \times 10^{4}$ W/m$^{2}$ — the flux the inner wall would transfer to ambient if the surface reached ambient temperature; minimizing $J$ therefore maximizes the heat rejected through the casing.  The Lagrange multiplier $\ell$ weights the material cost; it is a balance parameter of the optimisation, calibrated by the pipeline (a geometric scan locating the knee of the volume-fraction curve), not part of the problem statement. — the optimiser keeps only the material whose temperature reduction pays for itself.
The state $\left[ u\right] \in \mathbb P_{2}(D)$ is the solution of the linear problem
\begin{dmath*}
\int_{D} k \raisebox{\depth}{$\chi$} \left(  \nabla u  \cdot  \nabla \delta u  \right) \,dx + \int_{\Gamma_i} h_H u \delta u \,ds_x + \int_{\Gamma_o} h u \delta u \,ds_x = \int_{\Gamma_i} h_H u_{H} \delta u \,ds_x + \int_{\Gamma_o} h u_{amb} \delta u \,ds_x \qquad \forall\; {\delta u\in \mathbb P_{2}(D)} .
\end{dmath*}
The adjoint state $\left[ \tilde u\right]$ to the direct state $\left[ u\right]$ is the solution of
\begin{dmath*}
\int_{D} k \raisebox{\depth}{$\chi$} \left(  \nabla \delta u  \cdot  \nabla \tilde u  \right) \,dx + \int_{\Gamma_i} h_H \delta u \tilde u \,ds_x + \int_{\Gamma_o} h \delta u \tilde u \,ds_x = \int_{\Gamma_o} q_b \delta u \,ds_x \qquad \forall\; {\delta u\in \mathbb P_{2}(D)} ,
\end{dmath*}
where the adjoint load $\int_{\Gamma_o} q_b \delta u \,ds_x$ comes from the partial derivative $\partial_{u} J(\delta u) = - \int_{\Gamma_o} q_b \delta u \,ds_x$ of the outer-surface temperature term in the objective.
In order to compute the shape derivative, we consider a shape perturbation $\mathbf x = T_t(\mathbf X) = \mathbf X + t \mathbf  V(\mathbf X)$ with a suitable velocity field $\mathbf V$, vanishing on the fixed outer circle.  The perturbed Lagrangian $\mathcal G$ using the Lagrangian state $\left[ u\right]^{L}$ is given by
\begin{dmath*}
\mathcal G(t,\left[ u\right]^{L},\left[ \delta u\right]^{L}) = \int_{D} \ell \raisebox{\depth}{$\chi$}  \operatorname{det}\mathbf F( t )  \,dx + \int_{\Gamma_o} q_b \left( u_{amb} - u^{L}\right) \,ds_x +  \int_{D} k \raisebox{\depth}{$\chi$}  \operatorname{det}\mathbf F( t )  \left( \left(  \nabla u^{L} \cdot  \mathbf F^{-1}(t) \right) \cdot \left(  \nabla \delta u^{L} \cdot  \mathbf F^{-1}(t) \right) \right) \,dx + \int_{\Gamma_i} h_H \left( u^{L} - u_{H}\right) \delta u^{L} \,J_{s}(t) \,ds_x + \int_{\Gamma_o} h \left( u^{L} - u_{amb}\right) \delta u^{L} \,ds_x ,
\end{dmath*}
where $\mathbf F(t) = \mathbf I + t \nabla \mathbf V$ is the deformation gradient and $J_{s}(t)$ the boundary Jacobian of the moving inner boundary $\Gamma_i$; the fixed casing $\Gamma_o$ carries no Jacobian because $\mathbf V$ vanishes there.
For the direct state $\left[ u\right]$ and the adjoint state $\left[ \tilde u\right]$, the volume expression of the shape derivative, including the curvature term on the moving inner boundary, is given by
\begin{dmath*}
DJ(\Omega)(\mathbf V) = \int_{D} \left(\operatorname{div} {\mathbf V} \left(  \nabla u  \cdot  \nabla \tilde u  \right) - \left(  \nabla \tilde u  \cdot \left(  \nabla u  \cdot  \nabla {\mathbf V}  \right) \right) - \left(  \nabla u  \cdot \left(  \nabla \tilde u  \cdot  \nabla {\mathbf V}  \right) \right)\right) k \raisebox{\depth}{$\chi$} \,dx + \int_{D} \ell \raisebox{\depth}{$\chi$} \operatorname{div} {\mathbf V} \,dx + \int_{\Gamma_i} \kappa h_H \left( u - u_{H}\right) \tilde u \left(\mathbf V \cdot \mathbf n\right) \,ds_x ,
\end{dmath*}
where $\kappa$ is the curvature of $\Gamma_i$ and $\mathbf n$ its outward unit normal.  The inner circle $\Gamma_i$ is a design boundary — its shape is optimised — while the outer circle $\Gamma_o$ is a fixed boundary whose position is prescribed by the pipe geometry.
The initial design is a rotationally symmetric seed that covers both the heating wall and the casing while giving the optimiser enclosed interfaces to work from — a ring of evenly spaced holes on the mid-radius circle, one of the standard seeds of the pipeline's template library (the exact hole count and radius are numerical parameters chosen there, not part of the problem statement).
\end{document}
